\documentclass[10pt,a4paper]{amsart}
\usepackage[margin=19mm]{geometry}
\usepackage{amsmath,amssymb,mathtools}
\usepackage[scr=rsfs,cal=euler]{mathalfa}
\usepackage{xfrac}
\usepackage[unicode,hidelinks,bookmarks=false]{hyperref}
\newcommand{\ie}{i.e.\ }
\newcommand{\cf}{cf.\ }
\newcommand{\resp}{respectively\ }
\newcommand{\Subsection}{\subsection}
\providecommand{\nobreakdash}{\nobreak\hbox{-}}
\setlength{\emergencystretch}{2em}
\multlinegap0pt
\theoremstyle{plain}
\newtheorem{theorem}{Theorem}
\newtheorem{lemma}{Lemma}
\newtheorem{proposition}{Proposition}
\newtheorem{corollary}{Corollary}
\theoremstyle{definition}
\newtheorem{definition}{Definition}
\newtheorem{example}{Example}
\theoremstyle{remark}
\newtheorem{remark}{Remark}
\title[Nevanlinna theory of Hahn difference operators]{Nevanlinna theory of the Hahn difference operators and its applications: the second main theorem for Hahn difference operators}
\author{Ling Wang, Jianren Long, Xuxu Xiang}
\date{Source: arXiv:2509.18707v3 (CC BY 4.0)}
\begin{document}
\maketitle
\noindent\emph{Source and licence.} Nevanlinna theory of the Hahn difference operators and its applications. Version arXiv:2509.18707v3 (CC BY 4.0), distributed under the Creative Commons Attribution 4.0 International licence, \url{http://creativecommons.org/licenses/by/4.0/}, as recorded in the arXiv licence metadata for this exact version and shown on the abstract page \url{https://arxiv.org/abs/2509.18707v3}. Full licence notice: licenses/arxiv-2509.18707-CC-BY-4.0.txt.\par\smallskip

\noindent\textit{Editorial scope.} Counted unit: the article's Theorem with source label \texttt{thm2.5} (source0.tex lines 268--278), the second main theorem for the Hahn difference operator, stated as a defect relation for meromorphic functions of zero order together with the associated counting functions, delivered together with the article's complete original proof of it (source0.tex lines 279--318, one \texttt{proof} environment of 40 lines, adjacent to the statement). No mathematics is written, repaired or re-derived: statement and proof are the article's own text line for line, in the article's own notation (the Hahn difference operator, the Nevanlinna characteristic, the counting functions and the logarithmic density exceptional set), and no retained line is substituted, re-flowed or omitted. Required context retained uncounted: the statement of the article's earlier theorem (source0.tex lines 223--230) that the counted proof invokes. The proof of that context result is not delivered. Nothing else of the article is retained: its other theorems and its applications are omitted. No citation command occurs in any retained line, so the article's bibliography is not delivered and the wrapper carries no bibliography entry, although the article ships one. Every cross-reference inside the retained spans resolves inside the delivered document. Editorial formatting only: an AMS \texttt{amsart} preamble that restates the environments the retained lines use, namely the article's declarations of Theorem, Lemma, Proposition, Corollary, Definition, Example and Remark, and the macros that the retained lines need. The article is typeset with the Springer Nature \texttt{sn-jnl} class; that class file is present in the article's e-print but is neither loaded nor redistributed here, and the wrapper is an amsart document. Numbering differs from the article, because the article numbers its theorem-like environments continuously and only a subset is retained: the retained context theorem prints as Theorem 1 and the counted theorem as Theorem 2. No picture environment and no graphical inclusion occurs in any retained line, and the package delivers no image asset. One bundled source file, \texttt{sources/185537/source0.tex}, accompanies the delivered item as the exact source of every retained span. Every retained span is recorded in \texttt{delivery-evidence.json} with method \texttt{exact\_tex}.
\begin{theorem}\label{thm3.1}
	Suppose that $g$ is a non-constant meromorphic function with zero order. Then for any $0<|q|<1$,
	\begin{align*}
		m\left(r,\frac{\mathcal{D}_{q,c}g(z)}{g(z)}\right)=o(T(r,g)),
		m\left(r,\frac{\mathcal{D}_{q,c}^{k}g(z)}{g(z)}\right)=o(T(r,g))
	\end{align*}
	hold on a set of logarithmic density 1.
\end{theorem}

\begin{theorem}\label{thm2.5}
	Suppose that \(g\) is a transcendental meromorphic function with zero order and that
	\begin{align*}
		g^{n}(z)P(z,g)=Q(z,g),
	\end{align*}
	where \(P(z,g), Q(z,g)\) are Hahn difference polynomials in \(g\) and its any orders Hahn difference operators, and the degree of \(Q(z,g)\) is at most \(n\). Then
	\begin{align*}
		m(r,P(z,g))=o(T(r,g))
	\end{align*}
	as \(r\to\infty\).
\end{theorem}

\begin{proof}
	\indent We abbreviate \(P(z,g)\) as \(P(z)\) in our proof. According to the definition of proximity function \(m(r,g)\), then
	\begin{align*}
		m(r,P(z))&=\frac{1}{2\pi}\int_{0}^{2\pi}\log^{+}|P(re^{i\theta})|d\theta\\
		&=\frac{1}{2\pi}\int_{E}\log^{+}|P(re^{i\theta})|d\theta+\frac{1}{2\pi}\int_{E^{c}}\log^{+}|P(re^{i\theta})|d\theta,
	\end{align*}
	where set \(E=\{\theta\in[0,2\pi):|g(re^{i\theta})|<1\}\), and \(E^{c}\) is the complementary set of \(E\). On the set \(E\), denote
	\begin{align*}
		G(z)=a_{I}(z)g(z)^{l_{0}}(\mathcal{D}_{q,c}g(z))^{l_{1}}\cdots(\mathcal{D}^{v}_{q,c}g(z))^{l_{v}},
	\end{align*}
	then,
	\begin{align*}
		|G(z)|\leq|a_{I}|\left|\frac{\mathcal{D}_{q,c}g(z)}{g(z)}\right|^{l_{1}}\cdots\left|\frac{\mathcal{D}^{v}_{q,c}g(z)}{g(z)}\right|^{l_{v}}.
	\end{align*}
	Hence,
	\begin{align*}
		\frac{1}{2\pi}\int_{E}\log^{+}|G(re^{i\theta})|d\theta&\leq m(r,a_{I})+O\left(\sum_{t=1}^{v}m\left(r,\frac{\mathcal{D}^{v}_{q,c}g(z)}{g(z)}\right)\right)\\
		&+o(T(r,g)).
	\end{align*}
	Combining the Theorem \ref{thm3.1}, then
	\begin{align*}
		\frac{1}{2\pi}\int_{E}\log^{+}|P(re^{i\theta})|d\theta=\sum_{I}\int_{E}\log^{+}|G(re^{i\theta})|d\theta+O(1)=o(T(r,g)).
	\end{align*}
	\indent On the other hand, on the set \(E^{c}=\{\theta\in[0,2\pi):|g(re^{i\theta})|\geq1\}\). The Hahn difference polynomial \(Q(z,g)\) is the sum of terms of the form
	\begin{align*}
		a_{I}(z)g(z)^{l_{0}}(\mathcal{D}_{q,c}g(z))^{l_{1}}\cdots(\mathcal{D}^{v}_{q,c}g(z))^{l_{v}},~l_{0}+l_{1}+\cdots+l_{v}=I\leq n,
	\end{align*}
	\(v\) is a non-negative integer. So,
	\begin{align*}
		|P(z)|&=\left|\frac{1}{g^{n}(z)}\sum_{I}a_{I}(z)g(z)^{l_{0}}(\mathcal{D}_{q,c}g(z))^{l_{1}}\cdots(\mathcal{D}^{v}_{q,c}g(z))^{l_{v}}\right|\\
		&\leq\sum_{I}\left|a_{I}(z)\left|\frac{\mathcal{D}_{q,c}g(z)}{g(z)}\right|^{l_{1}}\cdots
		\left|\frac{\mathcal{D}^{v}_{q,c}g(z)}{g(z)}\right|^{l_{v}}\right|.
	\end{align*}
	Hence,
	\begin{align*}
		 \frac{1}{2\pi}\int_{E^{c}}\log^{+}|P(re^{i\theta})|d\theta&\leq\left(\sum_{t=1}^{v}m\left(r,\frac{\mathcal{D}^{t}_{q,c}g(z)}{g(z)}\right)+m(r,a(z))\right)\\
		&=o(T(r,g)).
	\end{align*}
	Therefore, Theorem \ref{thm2.5} is completely proved.
\end{proof}

\end{document}
